% ID: SEL-D0008-BD-E07
% ESTADO: cotejada. Segunda pasada visual completa y borrador compilado revisados.
% Fuente: Selectividad/matematicas_aplicadas/ORD_SUPLENTE_Modelo-4-Enunciados (Mat-Ap-2021).pdf
% SHA-256: 74b067b81be20aa0262a0cc58795f3b25488aa9c749d825b7b33f656dd1ae089
% Página: 2; bloque D; ejercicio original 7.

Se quiere estudiar la proporción de ciudadanos enfermos de COVID-19 en una
determinada población. Para ello, se elige una muestra al azar de 1000
ciudadanos, revelándose que el 15\% de ellos están enfermos.
\begin{enumerate}[label=\alph*)]
  \item \textbf{(1.5 puntos)} Calcule un intervalo de confianza al 95\%,
  para estimar la proporción real de enfermos de COVID-19 en dicha población.
  \item \textbf{(1 punto)} Determine el tamaño muestral mínimo para que,
  con el mismo nivel de confianza y la misma proporción muestral anteriores,
  el error que se cometa al estimar la proporción de ciudadanos enfermos de
  COVID-19 en esa población sea inferior al 1\%.
\end{enumerate}
